\section{Structural validation and archived measurements}
\label{sec:archived}

\paragraph{A structural input for compression.}
The 1,552-box family used in the validation has the following explicit
construction. For each of the six coordinate pairs in 4D, take 256
corners with pair extents $(t,257-t)$, $1\le t\le256$, and extent 300
in the other coordinates. Add 16 corners
$(161+i,176-i,161+j,176-j)$, where $0\le i<16$ and $j=3i\bmod16$.
The first 1,536 corners become pair constraints, leaving long arrays
with relatively few hard boxes. Shared extents are deliberate.

The current Numerical Chan driver executes six default-policy compressions
on this family and returns $6\,372\,330\,876$, agreeing with Chan $d/2$.
These events occur on different branches; they do not by themselves
test repeated compression along one path. The separate checkpoint
tests in Section~\ref{sec:software} cover that case.

An earlier recorded run on this same family illustrates the storage
effect mentioned in Section~\ref{sec:compression}: its first compression
changes one term into 296 terms and increases the stored-entry count
from 1,710 to 6,936. Compression shortens the coordinate histories but
need not immediately shrink the whole representation. These historical
event counts are not runtime measurements of the current driver.

\paragraph{Where to find the earlier experiments.}
The repository preserves the
\href{https://github.com/emmerichmtm/FastHVChan/blob/main/NumericalChanHVND/benchmarks/compiled_2d_10d/RESULTS.md}{earlier 2D--10D comparison}
at solver revision \texttt{38587e7}, including its small inputs, exact
oracle checks and compression probes. The larger experiment also retains
an \href{https://github.com/emmerichmtm/FastHVChan/blob/main/NumericalChanHVND/benchmarks/larger_4d_6d/RESULTS_ALL_VARIANTS.md}{all-variant detail report}.
Its extra implementation is an earlier compiled driver with a different
compression schedule. Those measurements remain available for
reproducibility; they are not included in this paper's main timing tables.

Raw inputs and results have not been changed or recomputed. The totals
in Section~\ref{sec:experiments} count only its three displayed methods;
the full archived experiment has additional rows and therefore different
aggregate counts. The current numerical method's complexity is justified
by the local policy and proof in Section~\ref{sec:complexity}, rather
than by historical timings or compilation alone.
